% Generated draft: requires mathematical and source review.
\documentclass[11pt]{article}
\usepackage[margin=25mm]{geometry}
\usepackage{fontspec}
\setmainfont{Latin Modern Roman}
\usepackage{amsmath,amssymb,mathtools,unicode-math}
\setmathfont{Latin Modern Math}
\usepackage{enumitem,needspace,xurl,hyperref}
\setlength{\parindent}{0pt}
\setlength{\parskip}{5pt}
\setlength{\emergencystretch}{3em}
\newcommand{\N}{\mathbb{N}}
\newcommand{\Z}{\mathbb{Z}}
\newcommand{\Q}{\mathbb{Q}}
\newcommand{\R}{\mathbb{R}}
\newcommand{\C}{\mathbb{C}}
\newcommand{\F}{\mathbb{F}}
\newcommand{\A}{\mathbb{A}}
\newcommand{\PP}{\mathbb{P}}
\newcommand{\E}{\mathbb{E}}
\newcommand{\eps}{\varepsilon}
\newcommand{\dd}{\,\mathrm{d}}
\begin{document}
\section*{UnsolvedMath ID: 20000387}
\textbf{Source problem code:} catalogue entry
\subsection{Definitions and mathematical statement}
Kiran Kedlaya's original question is: ``Look at principle (A) for tri-canonically embedded curves.'' The supplied verification says the full three-page source invokes Principles (A) and (B) but defines neither, so no exact original assertion can be recovered. The record's clean field gives contextual evidence rather than the principle: a question of Melanie Matchett Wood about i.i.d. local point-count models, neighboring fixed-field large-genus and complete-intersection questions, and earlier plane-curve/complete-intersection models.

The attached report adopts the following explicit reconstruction. Fix \(\F_q\) and \(g\geq2\). Embed each smooth projective geometrically connected genus-\(g\) curve by its complete \(3K_C\) system into \(\PP^r\), \(r=5g-6\), and choose uniformly among the projective framed images. It asserts that the induced measure on arithmetic curve isomorphism classes is proportional to \(1/|\operatorname{Aut}_{\F_q}(C)|\). Define the weighted factorial moments
\[
\mu_{g,s}=\frac{\sum_{[C]}(\#C(\F_q))_s/|\operatorname{Aut}_{\F_q}(C)|}
{\sum_{[C]}1/|\operatorname{Aut}_{\F_q}(C)|},
\]
where \((a)_s=a(a-1)\cdots(a-s+1)\), and let \(Q=\#\PP^r(\F_q)\). Write \(L_{r,s}\) for the number of ordered projectively independent ambient \(s\)-tuples, so
\[
L_{r,s}=\prod_{i=0}^{s-1}\left(Q-\frac{q^i-1}{q-1}\right).
\]
These definitions make explicit the normalization implicit in the source summary.

The report proposes that \(3K_C\) is sharply \((4g-6)\)-very ample, where \(k\)-very ampleness means surjectivity of restriction to every length-\(k+1\) zero-dimensional subscheme. Consequently, for \(1\leq s\leq4g-5\), its asserted incidence probability at a fixed ordered distinct ambient tuple is
\[
\Pr(P_1,\ldots,P_s\in C)=
\begin{cases}
\mu_{g,s}/L_{r,s}&\text{if the tuple is projectively independent},\\
0&\text{if it is dependent}.
\end{cases}
\]
For an independently chosen uniformly random ordered distinct ambient tuple, the averaged probability is instead \(\mu_{g,s}/(Q)_s\). With marginal \(p_g=\mu_{g,1}/Q\), the two proposed normalized ratios are
\[
\frac{\Pr(P_1,\ldots,P_s\in C)}{p_g^s}
=\frac{\mu_{g,s}}{\mu_{g,1}^s}\frac{Q^s}{L_{r,s}}
\quad\text{on independent tuples},
\]
\[
\frac{\Pr(\text{all points of a random distinct tuple lie on }C)}{p_g^s}
=\frac{\mu_{g,s}}{\mu_{g,1}^s}\frac{Q^s}{(Q)_s}.
\]
The report identifies their fixed-\(q\), fixed-\(s\), large-genus relative-factorization question with the intrinsic marked-moduli factorial-moment problem. The exact ensemble and limiting conventions need verification.

Collinear triples give zero joint incidence in this range, obstructing exact or uniform relative i.i.d. incidence. The report explicitly says this does not disprove absolute asymptotic factorization, since the marginal tends to zero. It leaves the intrinsic large-genus moment question open rather than inferring it from plane-curve or high-degree complete-intersection results.
\subsection{Short English statement}
The original principle is missing. Under an explicit framed-tricanonical reconstruction, the report gives exact incidence formulas, a dependent-tuple obstruction to relative independence, and a remaining factorial-moment question distinct from absolute asymptotic factorization.
\subsection{Sources}
\begin{itemize}
\item AIM workshop ``Arithmetic statistics over finite fields and function fields,'' Question 9, attributed to Kiran Kedlaya: \url{https://aimath.org/pastworkshops/arithstatffieldproblems.pdf}. The supplied verification says neither Principle (A) nor (B) is defined in the three-page source. No PDF was accessed here.
\item Local canonical registry \texttt{aim-algebraic-number-theory-notes.json}, \texttt{notes[38]}, tag \texttt{question}; permanent UnsolvedMath ID \(20000387\), source code \texttt{catalogue entry}. Original asks only about Principle (A) for tricanonical curves. The clean field instead gives contextual evidence; original remarks and literature field are empty.
\item Supplied contextual labels [AIM14b], for Melanie Matchett Wood's i.i.d.-model question; [BDFL10], for asymptotically independent conditions in a plane-curve model; and [BK12], for Bucur--Kedlaya's i.i.d. Bernoulli model for high-degree complete intersections. Nearby Question 8 concerns fixed-\(q\), large-genus counts and Question 10 complete intersections. No full bibliographic entries are supplied.
\item Supplied research attempt 1 proposing the exact incidence and relative-factorization formulas. It is machine generated, not peer reviewed, and requires independent verification before citation.
\end{itemize}
\end{document}
